\section{Multimodal Joint Analysis}

This section forms the core of the report. Rather than treating images and text as independent datasets, each subsection examines the dataset at the intersection of both modalities, revealing how visual content and textual descriptions relate to each other.

All statistics were computed directly from the actual RSITMD JSON annotation file using the dataset's built-in tokenization. Visual insights beyond metadata are inferred from caption content since pre-computed image features are not included in the dataset distribution.

\subsection{Image Category Distribution}

RSITMD contains 33 land use categories distributed unevenly across 4,291 training images. Understanding this distribution is critical because category imbalance propagates directly into the text modality: categories with more images contribute more captions, and vocabulary associated with frequent categories dominates overall word statistics.

The training set averages approximately 130 images per category, but the actual range spans from 6 (plane) to 227 (storagetanks) --- a \textbf{37.8:1 imbalance ratio}. Storage Tanks alone accounts for 5.3\% of all training images, while the smallest five categories combined (baseballfield, bareland, intersection, boat, plane) represent only 4.3\%.

\begin{figure}[H]
\centering
\begin{minipage}{\linewidth}
\resizebox{0.9\textwidth}{!}{\includegraphics{figures/ia_01_category_distribution_train.png}}
\end{minipage}
\caption{Category distribution across 33 land use categories in the training set. Bars are sorted in descending order; the imbalance between storagetanks (227) and plane (6) is clearly visible.}
\end{figure}

An important practical implication for multimodal retrieval: the dataset's skew means that image-text matching models will encounter majority-category pairs far more frequently during training. Without class-weighted sampling or data augmentation, minority categories (plane, boat, intersection) will receive insufficient training signal in both visual and textual embeddings.

The test set mirrors a similar but less extreme imbalance, with 32 categories (plane absent) averaging 14.1 images per category.

\subsection{Caption Characteristics}

Each of the 23,715 captions in RSITMD is a short, structured English description. Using the dataset's built-in tokenization (whitespace splitting, lowercase), caption word counts follow a near-normal distribution centered at 10 words, with mean \textbf{10.26} and standard deviation \textbf{3.73} in the training set. The minimum is 3 words, the maximum is 34. This tight clustering is characteristic of template-driven captioning.

\textbf{Vocabulary.} The training vocabulary contains 3,567 unique words (3,470 after stopword removal), which is relatively compact compared to general-purpose captioning datasets like COCO ($\sim$10,000 words). The test vocabulary of 1,218 unique words overlaps by 87.8\% with the training vocabulary (1,003 shared words out of 1,143 test clean words), meaning the test set uses almost exclusively words seen during training.

\begin{figure}[H]
\centering
\begin{minipage}{\linewidth}
\resizebox{0.88\textwidth}{!}{\includegraphics{figures/ta_02_word_frequency_train.png}}
\end{minipage}
\caption{Top 10 most frequent content words in training captions (stopwords removed). Color words (green, white) and object nouns (trees, buildings, road, river) dominate, confirming the visual nature of the captioning task.}
\end{figure}

\textbf{Word Frequency.} The most frequent content words across all captions are strongly visual: \textit{green} (4,567 occurrences), \textit{trees} (3,505), \textit{buildings} (2,921), \textit{white} (2,185), and \textit{road} (2,177). These are not abstract nouns but direct references to observable aerial features.

\textbf{N-gram Patterns.} The top bigrams are: \textit{green trees}, \textit{parking lot}, \textit{both sides}, \textit{many green}, \textit{dark green}, \textit{many buildings}, \textit{residential area}, \textit{football field}, \textit{several buildings}, and \textit{swimming pool}. These patterns reveal two dominant constructions: object-attribute pairs (green trees, dark green) and multi-object phrases (parking lot, many buildings).

\begin{figure}[H]
\centering
\begin{minipage}{\linewidth}
\resizebox{0.88\textwidth}{!}{\includegraphics{figures/ta_04_bigram_frequency_train.png}}
\end{minipage}
\caption{Top 10 bigrams in training captions (stopwords removed). Object-attribute pairs (green trees, dark green) and multi-object phrases (parking lot, many green) dominate, reflecting the template-driven annotation strategy.}
\end{figure}

The compact vocabulary and consistent length make RSITMD well-suited for fixed-length text encoding with token/padding masks. The dominance of concrete visual vocabulary (colors, objects, spatial relations) means text embeddings will capture primarily perceptual information rather than abstract semantic concepts.

\subsection{Visual Language in Captions}

The captions in RSITMD are fundamentally visual --- they describe what can be seen in aerial images rather than providing narrative or abstract interpretations. This section quantifies the visual vocabulary that forms the bridge between text and image modalities.

\textbf{Color Words.} Color is the most pervasive visual attribute in the captions. Approximately \textbf{41.5\%} of all training captions (8,904 out of 21,455) contain at least one color word. At the image level, 59.2\% of training images have at least one caption mentioning a color. The complete color word distribution is:

\begin{table}[H]
\centering
\caption{Color Word Frequency in Captions (Training Set)}
\label{tab:color_words}
\begin{tabular}{|l|c|c|l|}
\hline
\textbf{Color} & \textbf{Occurrences} & \textbf{\% All Caps} & \textbf{Semantic Relevance} \\
\hline
    green & 4,567 & 21.28\% & Vegetation, grass, trees \\
    white & 2,185 & 10.18\% & Buildings, sand, rooftops \\
    red & 1,001 & 4.67\% & Rooftops, sports fields \\
    blue & 939 & 4.38\% & Swimming pools, sea, sky \\
    gray & 796 & 3.71\% & Roads, urban structures \\
    yellow & 792 & 3.69\% & Rooftops, land markings \\
    grey & 502 & 2.34\% & Roads, urban areas \\
    brown & 447 & 2.08\% & Soil, bare land, wood \\
    black & 424 & 1.98\% & Roads, shadows \\
    orange & 164 & 0.76\% & Rooftops, vehicles \\
    silver & 99 & 0.46\% & Vehicles, rooftops \\
    purple & 32 & 0.15\% & Shadows, vehicles \\
    pink & 24 & 0.11\% & Rooftops \\
    golden & 17 & 0.08\% & Rooftops, sand \\
\hline
\end{tabular}
\end{table}

The dominance of \textit{green} (19.96\% of all captions, 36.8\% of color word occurrences) reflects the aerial perspective: vegetation is one of the most reliable and informative visual features in satellite imagery. Color words are the primary modality bridge --- they directly connect textual descriptions to pixel-level visual content.

\textbf{Spatial Relationships.} Spatial prepositions encode the top-down layout of aerial scenes: \textit{near}, \textit{surrounded}, \textit{next to}, \textit{around}, \textit{beside}, and \textit{in the middle of} are the most common. Approximately \textbf{31.3\%} of captions contain at least one spatial preposition, and at the image level, \textbf{52.9\%} of training images have at least one caption mentioning a spatial relationship.

\begin{table}[H]
\centering
\caption{Spatial Preposition Frequency in Captions (Training Set)}
\label{tab:spatial_words}
\begin{tabular}{|l|c|c|l|}
\hline
\textbf{Preposition} & \textbf{Count} & \textbf{\% Caps} & \textbf{Spatial Meaning} \\
\hline
near & 2,034 & 9.5\% & Proximity \\
surrounded & 1,707 & 8.0\% & Enclosure \\
next to & 1,090 & 5.1\% & Adjacency \\
around & 1,025 & 4.8\% & Surrounding \\
beside & 582 & 2.7\% & Side adjacency \\
in the middle of & 303 & 1.4\% & Central location \\
in the center of & 177 & 0.8\% & Central location \\
above & 39 & 0.2\% & Vertical relation \\
below & 29 & 0.1\% & Vertical relation \\
\hline
\end{tabular}
\end{table}

\textbf{Top Objects.} The most frequently mentioned content nouns in captions are: \textit{trees} (3,505), \textit{buildings} (2,921), \textit{road} (2,177), \textit{river} (1,699), \textit{building} (1,478), \textit{parking} (1,353), \textit{cars} (1,220), \textit{square} (1,210), \textit{area} (1,247), and \textit{houses} (1,134). These objects are precisely the visual features most distinguishable from an aerial perspective.

The convergence of color words, spatial prepositions, and concrete object nouns creates a rich visual-linguistic embedding space: text encodes primarily what is visually observable, and images contain the corresponding pixel-level evidence.

\subsection{Caption Variability and Sample Pairs}

One of the most distinctive features of RSITMD is its 5-caption-per-image design. Each satellite image is described by five independent captions that paraphrase each other rather than providing radically different perspectives. Understanding this variability is essential for retrieval and matching applications.

\textbf{Variability Quantification.} We measured caption length variability as the standard deviation of word counts across the five captions for each image. The mean std is \textbf{1.81 words} (maximum 9.93), confirming that captions for the same image cluster tightly around the 10-word target. This tight clustering confirms the annotation followed a structured guideline specifying target length.

The paraphrase-style diversity is valuable for cross-modal learning because it provides multiple positive image-text pairs for the same visual content. However, the five captions offer limited linguistic diversity --- all follow the same template of describing colors, objects, and spatial relationships in roughly 10 words.

\begin{figure}[H]
\centering
\begin{minipage}{\linewidth}
\resizebox{0.9\textwidth}{!}{\includegraphics{figures/mm_01_caption_variability_train.png}}
\end{minipage}
\caption{Caption variability analysis. Left: histogram of within-image caption length standard deviation (mean = 1.81). Right: mean variability by category (top 5). Categories with higher variability (airport, bridge, sparseresidential) suggest less template-driven captioning for those scene types.}
\end{figure}

\textbf{Sample Image-Caption Pairs.} The following figure displays three real satellite images from the RSITMD training set alongside their five captions:

\begin{figure}[H]
\centering
\resizebox{0.97\textwidth}{!}{\includegraphics{figures/mm_02_sample_pairs_train.png}}
\caption{Sample image-caption pairs from the RSITMD training set. Each row (A, B, C) shows one satellite image alongside its five independent captions. Color badges ([1] red, [2] blue, [3] green, [4] orange, [5] purple) mark corresponding captions across rows.}
\end{figure}

These examples reveal three patterns. First, all captions for the same image reference the same key objects (baseball field, playground, runway, park), confirming they describe the same visual content. Second, captions vary in syntactic construction: \textit{adjacent to} vs \textit{next to}, \textit{in the middle of} vs \textit{located in the middle of}. Third, color and spatial information is consistently present: \textit{green}, \textit{white}, \textit{red}, \textit{near}, \textit{middle of} appear across all captions, providing reliable cross-modal alignment signals.

The paraphrase diversity is sufficient for training image-text matching models but does not provide the full range of linguistic variation in real-world aerial scene descriptions. For more robust models, data augmentation techniques (e.g., back-translation) could artificially expand caption diversity beyond the structured annotation template.
